\chapter{Bắt Nạt Học Đường}
\markboth{Bắt Nạt Học Đường}{School Bullying}

\section{Bắt Nạt Không Phải Chuyện Đùa}
\begin{truyen}{Bắt Nạt Không Phải Chuyện Đùa}{Bullying Is Not a Joke}
\chuhoa{C}{ả lớp cười khi một bạn bị gọi bằng biệt danh gắn với ngoại hình.}
\textit{(The whole class laughed when one student was called a nickname tied to appearance.)}

Người nói bảo: ``Đùa thôi mà.''
\textit{(The speaker said: ``It is just a joke.'')}

Nhưng chỉ người bị gọi mới biết chữ đó bám theo mình cả ngày như thế nào.
\textit{(But only the person being called knew how that word clung to them all day.)}

Nhiều kiểu bắt nạt sống nhờ vào một câu nguỵ trang: đùa thôi.
\textit{(Many forms of bullying survive behind one disguise: just kidding.)}

Nó khiến người gây ra nhẹ tội trong mắt đám đông và khiến nạn nhân cảm thấy mình yếu đuối nếu phản ứng.
\textit{(It makes the aggressor look harmless to the crowd and makes the victim feel weak for reacting.)}

Sự thật thô nhưng thật là thế này: nếu thứ em làm khiến người khác co rúm, xấu hổ, và sợ xuất hiện, thì đó không phải trò đùa.
\textit{(The blunt truth is this: if what you do makes another person shrink, feel ashamed, and fear showing up, it is not humor.)}

Đó là bạo lực.
\textit{(It is violence.)}

\end{truyen}

\begin{baihoc}
Dạy học sinh gọi đúng tên sự việc. Khi gọi đúng tên, các em mới đủ can đảm dừng nó.
\textit{(Teach students to name the event correctly. Only then will they have the courage to stop it.)}
\end{baihoc}

\begin{ghinhoanh}
\textbf{Short English to remember:}

Teach students to name the event correctly.

Only then will they have the courage to stop it.

\end{ghinhoanh}

\begin{tuhocnhanh}
\textbf{Cau hoi tu hoc / Self-study questions}

1. Van de chinh la gi? \textit{What is the main problem?}

2. Cach xu ly tot hon la gi? \textit{What is a better action?}

3. Mot cau tieng Anh em muon nho la cau nao? \textit{Which English sentence do you want to remember?}
\end{tuhocnhanh}
\ngancach

\section{Người Đứng Nhìn Cũng Có Phần}
\begin{truyen}{Người Đứng Nhìn Cũng Có Phần}{The Bystander Has a Part Too}
\chuhoa{M}{ột bạn bị xé vở.}
\textit{(One student had a notebook torn up.)}

Ba mươi bạn khác nhìn thấy.
\textit{(Thirty other students saw it.)}

Không ai trực tiếp ra tay ngoài hai đứa.
\textit{(Only two acted directly.)}

Nhưng sự im lặng của ba mươi người còn lại chính là thứ khiến hai đứa kia thấy mình có quyền.
\textit{(But the silence of the other thirty was exactly what made the two feel entitled.)}

Nhiều học sinh tự an ủi: ``Mình đâu làm gì.'' Đúng, nhưng các em cũng đã không làm gì để chặn lại.
\textit{(Many students comfort themselves with: ``I did nothing.'' True, but they also did nothing to stop it.)}

Không phải ai cũng đủ sức lao vào đối đầu.
\textit{(Not everyone is strong enough to intervene openly.)}

Nhưng tối thiểu có thể báo giáo viên, ngồi cạnh nạn nhân, hoặc từ chối cười theo.
\textit{(But at minimum, one can tell a teacher, sit beside the victim, or refuse to laugh along.)}

\end{truyen}

\begin{baihoc}
Trong bắt nạt học đường, trung lập thường vô tình đứng về phía kẻ mạnh. Học sinh cần học rằng im lặng cũng là một lựa chọn đạo đức.
\textit{(In school bullying, neutrality often ends up serving the stronger side. Students must learn that silence is also a moral choice.)}
\end{baihoc}

\begin{ghinhoanh}
\textbf{Short English to remember:}

In school bullying, neutrality often ends up serving the stronger side.

Students must learn that silence is also a moral choice.

\end{ghinhoanh}

\begin{tuhocnhanh}
\textbf{Cau hoi tu hoc / Self-study questions}

1. Van de chinh la gi? \textit{What is the main problem?}

2. Cach xu ly tot hon la gi? \textit{What is a better action?}

3. Mot cau tieng Anh em muon nho la cau nao? \textit{Which English sentence do you want to remember?}
\end{tuhocnhanh}
\ngancach

\section{Nạn Nhân Không Cần Lời Khuyên Rẻ Tiền}
\begin{truyen}{Nạn Nhân Không Cần Lời Khuyên Rẻ Tiền}{Victims Do Not Need Cheap Advice}
\chuhoa{K}{hi một bạn kể mình bị bắt nạt, điều đầu tiên các bạn khác thường nói là: ``Kệ đi, đừng để ý.''}
\textit{(When someone admits being bullied, the first thing others often say is: ``Ignore it. Do not care.'')}

Nghe đơn giản, nhưng người đang bị cô lập đâu thể bật tắt cảm xúc như công tắc.
\textit{(It sounds simple, but a person being isolated cannot switch emotions on and off like a button.)}

Có bạn cần được tin, cần được ngồi cạnh, cần một người lớn can thiệp, chứ không cần thêm một câu khuyên thể hiện rằng nỗi đau của mình quá bé để ai phải bận tâm.
\textit{(Some students need to be believed, need company, need an adult to intervene, not another piece of advice implying that their pain is too small to matter.)}

Cách giúp tử tế nhất đôi khi rất ngắn: "Tao tin mày.
\textit{(The kindest help is sometimes very short: "I believe you.)}

Chuyện này không ổn.
\textit{(This is not okay.)}

Tao sẽ đi cùng mày báo cô."
\textit{(I will go with you to tell the teacher.")}

\end{truyen}

\begin{baihoc}
Hỗ trợ thật không phải là ném lời khuyên cho nhanh. Hỗ trợ thật là chia bớt gánh nặng và đưa nạn nhân về phía an toàn.
\textit{(Real support is not tossing out quick advice. Real support shares the burden and moves the victim toward safety.)}
\end{baihoc}

\begin{ghinhoanh}
\textbf{Short English to remember:}

Real support is not tossing out quick advice.

Real support shares the burden and moves the victim toward safety.

\end{ghinhoanh}

\begin{tuhocnhanh}
\textbf{Cau hoi tu hoc / Self-study questions}

1. Van de chinh la gi? \textit{What is the main problem?}

2. Cach xu ly tot hon la gi? \textit{What is a better action?}

3. Mot cau tieng Anh em muon nho la cau nao? \textit{Which English sentence do you want to remember?}
\end{tuhocnhanh}
\ngancach

\section{Kẻ Bắt Nạt Cũng Thường Là Một Đứa Trẻ Có Vấn Đề}
\begin{truyen}{Kẻ Bắt Nạt Cũng Thường Là Một Đứa Trẻ Có Vấn Đề}{A Bully Is Often Also a Child in Trouble}
\chuhoa{N}{ói thế không phải để bao che.}
\textit{(That does not excuse anything.)}

Kẻ bắt nạt phải chịu trách nhiệm.
\textit{(A bully must be held responsible.)}

Nhưng nếu chỉ phạt mà không hiểu gốc, nhà trường chỉ dập lửa ở ngọn.
\textit{(But if a school only punishes without understanding the root, it is only beating flames at the surface.)}

Có đứa bị bạo lực ở nhà nên mang bạo lực đến trường.
\textit{(Some children are harmed at home and bring harm to school.)}

Có đứa yếu cảm giác giá trị nên phải dìm người khác để thấy mình mạnh.
\textit{(Some feel so little worth that they push others down to feel strong.)}

Có đứa chỉ học được rằng ai yếu thì đáng bị nghiền.
\textit{(Some simply learned that the weak deserve to be crushed.)}

Xử lý bắt nạt phải vừa cứng vừa sâu: chặn hành vi ngay, bảo vệ nạn nhân ngay, nhưng cũng phải giáo dục lại kẻ gây ra trước khi nó lớn lên thành một người tàn nhẫn hơn.
\textit{(Handling bullying requires both firmness and depth: stop the behavior immediately, protect the victim immediately, and also re-educate the aggressor before that child grows into a crueler adult.)}

\end{truyen}

\begin{baihoc}
Kỷ luật tốt không chỉ để phạt. Kỷ luật tốt phải bảo vệ nạn nhân và chặn quá trình một đứa trẻ đang học cách làm ác.
\textit{(Good discipline is not only for punishment. It must protect the victim and interrupt a child who is learning how to do harm.)}
\end{baihoc}

\begin{ghinhoanh}
\textbf{Short English to remember:}

Good discipline is not only for punishment.

It must protect the victim and interrupt a child who is learning how to do harm.

\end{ghinhoanh}

\begin{tuhocnhanh}
\textbf{Cau hoi tu hoc / Self-study questions}

1. Van de chinh la gi? \textit{What is the main problem?}

2. Cach xu ly tot hon la gi? \textit{What is a better action?}

3. Mot cau tieng Anh em muon nho la cau nao? \textit{Which English sentence do you want to remember?}
\end{tuhocnhanh}
\ngancach

\section{Tên Gọi Dính Như Keo}
\begin{truyen}{Tên Gọi Dính Như Keo}{A Nickname That Sticks Like Glue}
\chuhoa{N}{gày xưa trong lớp học làng, một cậu bé nói lắp bị gọi bằng biệt danh suốt nhiều năm.}
\textit{(Long ago in a village classroom, a boy who stuttered was called by a mocking nickname for years.)}

Đám trẻ nói cho vui.
\textit{(The children said it for fun.)}

Thầy cũ nghe thấy cũng bỏ qua vì nghĩ chỉ là chuyện trẻ con.
\textit{(The old teacher heard it and ignored it, thinking it was only childish behavior.)}

Nhưng cậu bé dần ít nói hẳn, không dám đọc bài, rồi ghét luôn cả việc đến lớp.
\textit{(But the boy gradually stopped speaking, no longer dared to read aloud, and eventually hated coming to class.)}

Một lời chế giễu lặp lại đủ lâu sẽ thôi không còn là tiếng cười.
\textit{(A repeated mockery, carried on long enough, stops being laughter.)}

Nó biến thành căn phòng nhốt người ta lại.
\textit{(It becomes a room that imprisons a person.)}

\end{truyen}

\begin{baihoc}
Bắt nạt không cần đến nắm đấm mới gây thương tích. Nhiều vết thương học đường sống bằng lời nói và kéo dài rất lâu.
\textit{(Bullying does not require fists to injure. Many school wounds live through words and last a very long time.)}
\end{baihoc}

\begin{ghinhoanh}
\textbf{Short English to remember:}

Bullying does not require fists to injure.

Many school wounds live through words and last a very long time.

\end{ghinhoanh}

\begin{tuhocnhanh}
\textbf{Cau hoi tu hoc / Self-study questions}

1. Van de chinh la gi? \textit{What is the main problem?}

2. Cach xu ly tot hon la gi? \textit{What is a better action?}

3. Mot cau tieng Anh em muon nho la cau nao? \textit{Which English sentence do you want to remember?}
\end{tuhocnhanh}
\ngancach

\section{Thầy Giáo Mới Và Chiếc Ghế Cuối Lớp}
\begin{truyen}{Thầy Giáo Mới Và Chiếc Ghế Cuối Lớp}{The New Teacher and the Back-Row Seat}
\chuhoa{M}{ột thầy giáo mới vào lớp và thấy có một học sinh luôn bị đẩy xuống ngồi cuối, đồ dùng hay bị giấu mất.}
\textit{(A new teacher entered a class and noticed one student was always pushed to the back and often had belongings hidden.)}

Thầy không mắng ngay giữa lớp.
\textit{(He did not scold the class immediately.)}

Thầy đổi chỗ ngồi, quan sát thêm, rồi gọi riêng từng em lên hỏi.
\textit{(He changed the seating, observed more, and then called students in one by one.)}

Hóa ra cả lớp đã quen với việc một bạn bị đùa dai đến mức không ai còn thấy đó là chuyện lạ.
\textit{(It turned out the whole class had become so used to one student being targeted that no one saw it as unusual anymore.)}

Cái nguy của bạo lực lặp đi lặp lại là nó được bình thường hóa.
\textit{(The danger of repeated harm is that it becomes normalized.)}

\end{truyen}

\begin{baihoc}
Muốn chặn bắt nạt, người lớn phải nhìn được những điều lớp học đã coi là bình thường từ lâu.
\textit{(To stop bullying, adults must see what the classroom has long begun treating as normal.)}
\end{baihoc}

\begin{ghinhoanh}
\textbf{Short English to remember:}

To stop bullying, adults must see what the classroom has long begun treating as normal.

\end{ghinhoanh}

\begin{tuhocnhanh}
\textbf{Cau hoi tu hoc / Self-study questions}

1. Van de chinh la gi? \textit{What is the main problem?}

2. Cach xu ly tot hon la gi? \textit{What is a better action?}

3. Mot cau tieng Anh em muon nho la cau nao? \textit{Which English sentence do you want to remember?}
\end{tuhocnhanh}
\ngancach

\section{Kẻ Mạnh Miệng Nhưng Yếu Người}
\begin{truyen}{Kẻ Mạnh Miệng Nhưng Yếu Người}{Loud but Weak Inside}
\chuhoa{M}{ột cậu học sinh hay chặn đường bắt nạt bạn nhỏ hơn.}
\textit{(A student often cornered and bullied smaller classmates.)}

Trước đám đông em rất lì.
\textit{(In front of crowds he seemed fearless.)}

Nhưng khi cô tư vấn trò chuyện riêng, em bật ra chuyện mình thường xuyên bị chửi ở nhà và chưa bao giờ thấy mình có giá trị.
\textit{(But in a private counseling session, he revealed that he was often verbally abused at home and had never felt valuable.)}

Điều đó không làm hành vi của em bớt sai.
\textit{(This did not make his actions less wrong.)}

Nhưng nó cho thấy sự hung hăng đôi khi là chiếc áo khoác của yếu đuối.
\textit{(But it showed that aggression is sometimes the coat worn by weakness.)}

Nếu nhà trường chỉ ghét kẻ bắt nạt mà không hiểu cơ chế tạo ra nó, vòng lặp sẽ còn tiếp tục.
\textit{(If a school only hates the bully and does not understand the mechanism producing him, the cycle will continue.)}

\end{truyen}

\begin{baihoc}
Kỷ luật cần đi cùng hiểu biết. Không phải để mềm tay, mà để chặn đúng gốc của cái ác nhỏ đang lớn lên.
\textit{(Discipline must be paired with understanding. Not to go soft, but to interrupt the root of a small cruelty that is growing.)}
\end{baihoc}

\begin{ghinhoanh}
\textbf{Short English to remember:}

Discipline must be paired with understanding.

Not to go soft, but to interrupt the root of a small cruelty that is growing.

\end{ghinhoanh}

\begin{tuhocnhanh}
\textbf{Cau hoi tu hoc / Self-study questions}

1. Van de chinh la gi? \textit{What is the main problem?}

2. Cach xu ly tot hon la gi? \textit{What is a better action?}

3. Mot cau tieng Anh em muon nho la cau nao? \textit{Which English sentence do you want to remember?}
\end{tuhocnhanh}
\ngancach

\section{Tiếng Cười Trong Hành Lang}
\begin{truyen}{Tiếng Cười Trong Hành Lang}{The Laughter in the Corridor}
\chuhoa{M}{ột bạn nữ đi ngang hành lang và nghe sau lưng mình có tiếng cười.}
\textit{(A girl walked down the corridor and heard laughter behind her.)}

Không ai gọi tên em, nhưng em biết họ đang cười mình.
\textit{(No one called her name, yet she knew they were laughing at her.)}

Kiểu bắt nạt này khó nắm hơn vì nó để lại rất ít bằng chứng.
\textit{(This kind of bullying is harder to catch because it leaves little evidence.)}

Nhưng nạn nhân vẫn đau thật.
\textit{(But the victim still hurts for real.)}

Có những tổn thương học đường đến từ bầu không khí chứ không chỉ từ một hành vi cụ thể.
\textit{(Some school injuries come from an atmosphere, not only from one concrete act.)}

Một môi trường khiến học sinh phải đoán xem tiếng cười nào là dành cho mình là một môi trường đang có vấn đề.
\textit{(An environment where students must guess which laugh is directed at them is an environment with a problem.)}

\end{truyen}

\begin{baihoc}
Không phải bạo lực nào cũng rõ ràng. Người lớn cần tinh hơn với những dạng tổn thương không có dấu vết vật lý.
\textit{(Not every form of harm is obvious. Adults must become sharper in noticing injuries without physical marks.)}
\end{baihoc}

\begin{ghinhoanh}
\textbf{Short English to remember:}

Not every form of harm is obvious.

Adults must become sharper in noticing injuries without physical marks.

\end{ghinhoanh}

\begin{tuhocnhanh}
\textbf{Cau hoi tu hoc / Self-study questions}

1. Van de chinh la gi? \textit{What is the main problem?}

2. Cach xu ly tot hon la gi? \textit{What is a better action?}

3. Mot cau tieng Anh em muon nho la cau nao? \textit{Which English sentence do you want to remember?}
\end{tuhocnhanh}
\ngancach

\section{Bảo Vệ Bạn Làm Em Mất Vai Trò “An Toàn”}
\begin{truyen}{Bảo Vệ Bạn Làm Em Mất Vai Trò “An Toàn”}{Protecting Someone Costs the Safe Position}
\chuhoa{M}{ột học sinh chứng kiến bạn bị cô lập.}
\textit{(A student watched a classmate being isolated.)}

Em biết nếu mình bênh, nhóm kia có thể quay sang loại cả mình.
\textit{(She knew that if she defended him, the group might turn and exclude her too.)}

Im lặng là lựa chọn an toàn nhất.
\textit{(Silence was the safest option.)}

Và cũng là lựa chọn khiến em khó nhìn thẳng vào mình nhất.
\textit{(It was also the one that made it hardest for her to look herself in the eye.)}

Cuối cùng em chỉ làm một việc nhỏ: ngồi cạnh bạn ấy trong giờ ăn trưa.
\textit{(In the end she did one small thing: she sat beside him at lunch.)}

Nhỏ, nhưng đủ để phá nứt bức tường cô lập.
\textit{(Small, but enough to crack the wall of isolation.)}

Chống bắt nạt không phải lúc nào cũng là bài diễn văn lớn.
\textit{(Fighting bullying is not always a grand speech.)}

Nhiều khi nó chỉ là không bỏ một người ngồi một mình.
\textit{(Often it is simply refusing to leave someone alone.)}

\end{truyen}

\begin{baihoc}
Học sinh cần những hành động can đảm vừa tầm, không phải lúc nào cũng anh hùng kiểu sân khấu.
\textit{(Students need forms of courage they can realistically enact, not only heroic stage-like gestures.)}
\end{baihoc}

\begin{ghinhoanh}
\textbf{Short English to remember:}

Students need forms of courage they can realistically enact, not only heroic stage-like gestures.

\end{ghinhoanh}

\begin{tuhocnhanh}
\textbf{Cau hoi tu hoc / Self-study questions}

1. Van de chinh la gi? \textit{What is the main problem?}

2. Cach xu ly tot hon la gi? \textit{What is a better action?}

3. Mot cau tieng Anh em muon nho la cau nao? \textit{Which English sentence do you want to remember?}
\end{tuhocnhanh}
\ngancach

\section{Kỷ Luật Để Giữ Người, Không Chỉ Để Giữ Mặt Mũi Trường}
\begin{truyen}{Kỷ Luật Để Giữ Người, Không Chỉ Để Giữ Mặt Mũi Trường}{Discipline to Protect People, Not Just the School’s Image}
\chuhoa{C}{ó trường xử lý bắt nạt rất chậm vì sợ ồn ào, sợ phụ huynh bàn tán, sợ ảnh hưởng danh tiếng.}
\textit{(Some schools handle bullying slowly because they fear noise, parental talk, and damage to reputation.)}

Trong lúc người lớn lo giữ mặt mũi, học sinh bị hại vẫn phải bước vào lớp mỗi ngày với nỗi sợ nguyên vẹn.
\textit{(While adults worry about image, the harmed student still walks into class every day with intact fear.)}

Một ngôi trường có danh tiếng đẹp nhưng để học sinh sống trong sợ hãi thì cái đẹp đó rất rỗng.
\textit{(A school with a polished reputation but fearful students holds a hollow kind of beauty.)}

Kỷ luật đúng không phải để sự việc biến mất trên giấy tờ.
\textit{(Proper discipline is not to make incidents disappear from paperwork.)}

Kỷ luật đúng là để người yếu không bị bỏ mặc.
\textit{(Proper discipline is to ensure the vulnerable are not abandoned.)}

\end{truyen}

\begin{baihoc}
Danh tiếng nhà trường chỉ đáng giữ khi nó được xây trên sự an toàn thật của học sinh.
\textit{(A school’s reputation is worth preserving only when it rests on real student safety.)}
\end{baihoc}

\begin{ghinhoanh}
\textbf{Short English to remember:}

A school’s reputation is worth preserving only when it rests on real student safety.

\end{ghinhoanh}

\begin{tuhocnhanh}
\textbf{Cau hoi tu hoc / Self-study questions}

1. Van de chinh la gi? \textit{What is the main problem?}

2. Cach xu ly tot hon la gi? \textit{What is a better action?}

3. Mot cau tieng Anh em muon nho la cau nao? \textit{Which English sentence do you want to remember?}
\end{tuhocnhanh}
